\subsection{Свойства собственных значений компактного оператора}

\begin{definition}[Собственное значение]
	Пусть \(E\) --- комплексное банахово пространство и
	\[
	A\in\mathcal{L}(E).
	\]
	Число \(\lambda\in\mathbb{C}\) называется \textbf{собственным значением} оператора \(A\), если
	существует ненулевой вектор \(x\in E\), такой что
	\[
	Ax=\lambda x.
	\]
	Вектор \(x\) называется \textbf{собственным вектором}.
\end{definition}

\begin{remark}[Почему ноль особый]
	Для компактного оператора ненулевые собственные значения ведут себя почти как в конечномерной
	линейной алгебре: вне любой окрестности нуля их может быть только конечное число.
	
	Но \(0\) может быть собственным значением бесконечной кратности.
\end{remark}

\begin{example}
	Пусть оператор \(A:l_2\to l_2\) задаётся диагональной матрицей
	\[
	A=
	\begin{pmatrix}
		1 & 0 & 0 & \cdots\\
		0 & \frac12 & 0 & \cdots\\
		0 & 0 & \frac13 & \cdots\\
		\vdots & \vdots & \vdots & \ddots
	\end{pmatrix}.
	\]
	То есть
	\[
	A(x_1,x_2,x_3,\ldots)
	=
	\left(x_1,\frac{x_2}{2},\frac{x_3}{3},\ldots\right).
	\]
	Этот оператор компактен.
	
	Его ненулевые собственные значения:
	\[
	1,\frac12,\frac13,\ldots
	\]
	и они стремятся к нулю. При этом \(0\) является собственным значением бесконечной кратности
	для подходящих компактных операторов, поэтому \(0\) нельзя рассматривать так же, как
	ненулевые собственные значения.
\end{example}

\begin{theorem}[Собственные значения компактного самосопряженного оператора вне нуля]
	Пусть \(H\) --- гильбертово пространство и
	\[
	A\in\mathcal K(H)
	\]
	--- самосопряженный оператор. Тогда для любого
	\[
	\delta>0
	\]
	существует лишь конечное число собственных значений \(\lambda\), для которых
	\[
	|\lambda|\geq\delta.
	\]
\end{theorem}

\begin{proof}
	\textbf{Хотим получить противоречие с компактностью \(A\).}
	
	Предположим, что существуют различные собственные значения
	\[
	\lambda_1,\lambda_2,\ldots,
	\qquad
	|\lambda_n|\geq\delta.
	\]
	Выберем соответствующие нормированные собственные векторы
	\[
	Ae_n=\lambda_n e_n,
	\qquad
	\|e_n\|=1.
	\]
	
	\textbf{Хотим показать, что образы \(Ae_n\) не могут иметь сходящейся подпоследовательности.}
	
	Так как \(A\) самосопряжен, собственные векторы различных собственных значений
	ортогональны:
	\[
	e_n\perp e_m,
	\qquad n\neq m.
	\]
	
	Поэтому
	\[
	\begin{aligned}
		\|Ae_n-Ae_m\|^2
		&=
		\|\lambda_n e_n-\lambda_m e_m\|^2
		\\
		&=
		|\lambda_n|^2+|\lambda_m|^2
		\\
		&\geq
		2\delta^2.
	\end{aligned}
	\]
	Следовательно,
	\[
	\|Ae_n-Ae_m\|
	\geq
	\sqrt2\,\delta,
	\qquad n\neq m.
	\]
	
	Значит, последовательность
	\[
	\{Ae_n\}
	\]
	не содержит фундаментальной, а значит и сходящейся подпоследовательности.
	
	Но
	\[
	\|e_n\|=1,
	\]
	то есть \(\{e_n\}\) ограничена. Так как \(A\) компактен,
	последовательность \(\{Ae_n\}\) должна иметь сходящуюся подпоследовательность.
	
	Получили противоречие.
	
	Следовательно, для любого \(\delta>0\) существует лишь конечное число
	собственных значений, удовлетворяющих
	\[
	|\lambda|\geq\delta.
	\]
\end{proof}

\begin{remark}[Общий случай]
	В общем банаховом случае это утверждение доказывается похожей идеей, но вместо
	ортогональности используют теорему Рисса о почти перпендикуляре.
	
	На экзамене обычно достаточно понимать смысл: если бы ненулевых собственных значений
	вдали от нуля было бесконечно много, компактность дала бы сходящуюся подпоследовательность
	образов, но собственные векторы можно выбрать так, что образы остаются отделёнными друг от друга.
\end{remark}
